\documentclass[12pt,a4paper,oneside]{article}
\usepackage[brazil]{babel}
\usepackage[utf8]{inputenc}
\usepackage{olymp}
\usepackage{color}
\usepackage[dvipsnames]{xcolor}
\usepackage{listings}

\setcounter{problem}{4}

\begin{document}

\begin{problem}{Distintos}{distintos}{2}
 
Dada uma lista de $N$ números inteiros, seu programa deve determinar quantos números distintos existem na lista.

%\Notes
%
%Observações, se tiver.

\InputFile

A entrada contém duas linhas: a primeira contém um valor inteiro $N$, indicando a quantidade de números na lista; a segunda contém os $N$ números da lista, que são números inteiros $V_i$ separados por espaço em branco.

Restrições: $1 \leq N \leq 1000$, $-2^{31} \leq V_i \leq 2^{31}-1$ (ou seja, qualquer valor \texttt{int}).



\OutputFile

Seu programa deve gerar apenas uma linha de saída, informando a quantidade de números distintos da lista.

%\Notes
%Nesta questão, seu programa deve usar a função \texttt{fatorial} dada %acima exatamente como está, não a modifique! Sua
%tarefa é apenas criar o programa com a função \texttt{main}.

\Example

\begin{example}
\exmp{
10
1 2 3 4 5 6 7 8 9 10
}{
10
}%
\end{example}

\begin{example}
\exmp{
10
5 1 0 -2 0 5 1 0 -2 3
}{
5
}%
\end{example}

\begin{example}
\exmp{
3
0 0 0
}{
1
}%
\end{example}

\begin{example}
\exmp{
7
0 1 1 1 0 1 0
}{
2
}%
\end{example}

%Comentários sobre o exemplo, se necessário.

\end{problem}

\end{document}
